\documentclass[11pt]{article}
\usepackage[margin=1in]{geometry}
\usepackage{amsmath}
\usepackage{booktabs}
\usepackage{url}
\setlength{\emergencystretch}{2em}
\title{Tracker parameter noise and ablations of three controller choices}
\date{Prepared October 5, 2026}
\begin{document}
\maketitle

\section{Questions}
With the tracker's own target acceleration $A$ and sideslip $\beta$
published (commit \texttt{d1c31d6},
\path{ESTIMATOR_DIRECT_TARGET_20261005.tex}), two of 84 noisy encounters
end without a solved command. The user approved one remedy: add a bound on
the noise of the tracker's acceleration state to the semidefinite program
that selects its gains. The user also asked whether three controller choices
of \texttt{06af0a2} are needed and what they do:
\begin{itemize}
\item road recovery in the terminal set;
\item the control Lyapunov function (CLF) row on the estimate;
\item soft untightened collision rows in the slack prefix.
\end{itemize}
This report records the gain study, a prototype of a different estimator
remedy, and one single-change ablation per controller choice. No project
source file was changed.

\section{Tracker gains with an acceleration-noise objective}
\paragraph{Where the noise comes from.} The tracker state is $[\rho;q;s]$, and
$A=q^\top s/|q|$ and $\beta$ are reconstructed from the acceleration state
$s$. Radar noise is uniform within $\pm0.0283$~m per axis (standard
deviation 0.0163~m) at 80~Hz. Through the current gains $(17.12, 144.96,
558.93)$ its stationary per-axis effect is 0.0068~m in position,
0.054~m/s in velocity and 0.198~m/s$^2$ in $s$. The measured $A$ error has
a per-run RMS median of 0.27~m/s$^2$ between 0.5 and 3~s, 0.23 between 3 and
6~s, and 0.19 between 6 and 10~s. The radar noise therefore accounts for it.

\paragraph{What fixes the gain level.} The certificate requires the
Lipschitz dissipation inequality at the decay floor $1/T=0.721$~s$^{-1}$.
Here $T=1.387$~s is the domain transit time, and the Lipschitz constants are
10.36 and 2.01. For the current shape, the inequality holds at bandwidth
4.55 but not at 4.50; the design uses 4.70
(\path{TRACKER_NOISE_AND_DESIGN_ABLATIONS_20261005/certified-search.csv}).
Noise in $s$ grows roughly as the bandwidth to the power 2.5, so this floor
sets it.

\paragraph{Weighted objective.} The shape program bounds the stationary
error variance through a Gramian inequality that shares one Lyapunov
matrix with the dissipation inequality. That variance was weighted on the
velocity and acceleration errors as listed below. The table gives the
predicted per-axis noise after the existing bandwidth search
(\path{TRACKER_NOISE_AND_DESIGN_ABLATIONS_20261005/shape-weights.csv}).
\begin{center}
\begin{tabular}{lrrrr}
\toprule
weights (pos., vel., acc.) & bandwidth & vel.\ (m/s) & acc.\ (m/s$^2$) & certified acc.\ bound\\
\midrule
0, 1, 0 (current) & 4.70 & 0.0539 & 0.1981 & 121.4\\
0, 1, 0.01 & 4.84 & 0.0545 & 0.1982 & 122.5\\
0, 1, 0.1 & 5.22 & 0.0559 & 0.2014 & 127.5\\
0, 1, 1 & 5.29 & 0.0561 & 0.2025 & 128.9\\
0, 0, 1 & 5.30 & 0.0562 & 0.2026 & 129.0\\
\bottomrule
\end{tabular}
\end{center}
Weighting the acceleration error raised the true acceleration noise by up to
2\%. The program minimizes an upper bound shared with the certificate, and
that bound pulled the bandwidth up. A direct search over the three gains was
run next. Its objective was the exact stationary variance (Lyapunov
equation). Each candidate was checked with the code's own scaled
free-metric certificate. The lowest certified acceleration noise it found
was 0.1675~m/s$^2$, 15\% below the current gains, at gains $(12.69, 105.57,
389.35)$. The window fit that was removed had an $A$ error of 0.019. The
approved change was therefore not made: inside the certificate, the gains
cannot remove the noise that caused the failures.

\section{Prototype: a constant-parameter stage}
The noise in $s$ is high-frequency, since $s$ is the second derivative of the
radar position. $A$ and $\beta$ are constant by contract, so a first-order
stage behind the tracker can average this noise. The stage tracks the
tracker's reconstruction at the existing decay floor (time constant 1.39~s).
It is one more feedforward stage of the cascade, and its error is ISS with
unit gain from the tracker's parameter error. For that rate, the Lyapunov
equation of the tracker and stage predicts an acceleration noise of
0.0118~m/s$^2$, against 0.198 unfiltered
(\path{TRACKER_NOISE_AND_DESIGN_ABLATIONS_20261005/parameter-stage-noise.csv}).

Two versions were run on the 84 noisy encounters. v1 starts from the
tracker's value and filters at the floor rate. v2 starts as the running mean
of the tracker's values (gain $1/n$) until that gain reaches the floor
(\path{patches/param.patch}, \path{patches/param2.patch}). The published radii
are widened by the stage--tracker distance, as they were for the window fit.
Results are in Section~4. Both recovered 82 of 84, with different failures
from the tracker output:
\begin{itemize}
\item \emph{Seed 1, 15-m/s crossing (v1 and v2), 2.80~s.} The ego was in a
  sharp recovery (lateral velocity 1.86~m/s, yaw rate $-0.53$~rad/s), with
  the target 31~m away. The published $A$ was $-0.03$ (truth 0). Clarabel
  certified the primary problem infeasible, for the shifted plan and the
  fresh restart alike. There was still no solution with any single
  constraint group removed. The parameter estimate is not its cause; the
  trajectory differs.
\item \emph{v1, seed 5, 15-m/s braking lead, 0.75~s.} The stage lagged: $-0.15$
  against a true $-0.50$, while the tracker alone had reached $-0.68$ at
  0.3~s. The horizon grew to 173 holds, and the CLF stage stopped without a
  numerical result. v2 recovered in this run.
\item \emph{v2, seed 20261003, 15-m/s braking lead, 0.65~s.} The tracker's
  initial transient was positive, so the running mean was $+0.18$ at 0.35~s
  and $-0.17$ at 0.6~s. The horizon grew to 176 holds. Clarabel then stopped the CLF stage
  with insufficient progress (primal residual $1.2\times10^{-5}$).
\end{itemize}
The stage is not adopted. Its start-up needs a design, and tuning it further
on the same seeds would fit the evaluation set.

\section{Controller ablations}
Each variant changes one choice of \texttt{d1c31d6} back to its form before
\texttt{06af0a2} (\path{patches/road.patch}, \path{clf.patch},
\path{prefix.patch}):
\begin{itemize}
\item \emph{road}: no road-recovery rows, and the rollout stops on separation
  alone;
\item \emph{clf}: the CLF decrease is required over the whole ego enclosure;
\item \emph{prefix}: the untightened collision rows are hard at every node,
  including the prefix.
\end{itemize}
The variants ran as exported copies of the commit. They used the same 84 noisy
runs (seeds 20261003 and 1--5) and, for \emph{road}, the 14 exact runs. Ten
processes ran in parallel. Four runs that hit the 5-s solver time limit were
rerun without contention, and their rerun outcomes are counted:
\begin{itemize}
\item \emph{clf}: one run recovered and one still had no solution;
\item \emph{prefix}: the rerun recovered;
\item \emph{v1}: the rerun recovered.
\end{itemize}
\begin{center}
\begin{tabular}{lrrrrrr}
\toprule
noisy, 84 runs & recovered & no solution & not in 100 s & min.\ gap & min.\ road & $A$ RMS\\
\midrule
\texttt{d1c31d6} & 82 & 2 & 0 & 0.362 & 0.23 & 0.269\\
road recovery removed & 72 & 12 & 0 & 0.362 & 0.16 & 0.272\\
worst-case CLF & 38 & 9 & 37 & 0.356 & 0.11 & 0.276\\
hard prefix rows & 79 & 5 & 0 & 0.432 & 0.23 & 0.269\\
parameter stage v1 & 82 & 2 & 0 & 0.412 & 0.96 & 0.065\\
parameter stage v2 & 82 & 2 & 0 & 0.425 & 0.74 & 0.109\\
\bottomrule
\end{tabular}
\end{center}
No run collided or left the road. Gaps and road margins are in metres. The
$A$ RMS is the per-run median over 0.5--3~s, in m/s$^2$; between 3 and 6~s it
was 0.228 for \texttt{d1c31d6}, 0.032 for v1 and 0.030 for v2. With exact
states, removing road recovery left all 14 runs recovering. The other variants
do not act without an enclosure, so they were not run with exact states.
Per-run data: \path{TRACKER_NOISE_AND_DESIGN_ABLATIONS_20261005/ablation-metrics.csv}.

\paragraph{Road recovery.} Ten of the twelve failures are new; the other two
are the failures of \texttt{d1c31d6}. Nine of the new ones are 8-m/s
head-on-type encounters (head-on, accelerating head-on, curved head-on) that
failed between 1.70 and 2.25~s. In the last solved frame of each, the ego
headed 30--61~degrees into the left edge. It moved toward that edge at
3.2--7.7~m/s with 0.3--2.1~m left; these values come from the recorded
path-relative pose. Stopping that motion at full friction would
have needed more lateral room than remained (for example 2.00 against
1.31~m), so the road rows had no solution. With road recovery, the same nine
runs peaked at 53--59~degrees and 6.4--6.8~m/s toward the edge, and all
recovered with 0.6--1.3~m of margin. The tenth new failure (15-m/s curved
head-on, 55.05~s) was not near an edge.

\paragraph{CLF on the estimate.} With the worst-case CLF, 37 of the 42
15-m/s runs did not recover within 100~s, and 9 runs had no solution. The
braking ratio chattered. The mean absolute change between holds was 0.0353
(\texttt{d1c31d6}: 0.0106), and its standard deviation after 4~s was 0.119
(\texttt{d1c31d6}: 0.0011). This is the behaviour described for the original
design: near the path, the worst case cannot decrease below the enclosure's
own size, so the CLF stage minimizes it greedily.

\paragraph{Prefix rows.} Hard prefix rows added three failures, all in the
15-m/s accelerating head-on (seeds 20261003, 3 and 5), at 0.45--0.55~s near
the closest approach. The 8-m/s braking lead of seed 3, which fails at
2.65~s in \texttt{d1c31d6}, failed at 1.05~s. Nothing else differs between
the variants, so these runs fail because of the hard rows. In
\texttt{d1c31d6} the prefix soft rows carried slack above 1~mm in 1.68\% of
the noisy holds.

\paragraph{Horizon length.} The failures of \texttt{d1c31d6} and of both
parameter-stage versions that depend on $A$ occurred once the horizon had
grown to 84--176 holds. The CLF stage's numerical stalls occurred at
horizons of 173 and 176 holds.

\section{Scope}
\begin{itemize}
\item The closed loop is sensitive to small changes. A single change can move
  marginal runs either way, which is why the failure cases are listed
  individually.
\item The noise predictions are continuous-time stationary values for a
  white-noise equivalent of the sampled radar noise. They ignore the
  nonlinear terms and the ego estimation errors.
\item The direct search is local (fminsearch from five starts), so 0.1675 is
  the best certified value found, not a proven minimum.
\item The parameter stage is a prototype; its error bound beyond the
  unit-gain ISS statement was not developed.
\end{itemize}

\section{Reproduction}
The files below are in the folder
\path{TRACKER_NOISE_AND_DESIGN_ABLATIONS_20261005}.
\begin{itemize}
\item Study tables: \path{scripts/sdpTables.m.txt}; the weighted shape
  program is \path{patches/shape-weights.patch}.
\item Campaigns: \path{scripts/campaignOneRoot.m.txt}, run by
  \path{scripts/launchRoot.sh.txt}. Each variant root is commit
  \texttt{d1c31d6} exported with \texttt{git archive}, plus one patch from
  \path{patches/} and a link to \path{solver/}.
\item Summaries: \path{scripts/ablation.py.txt}.
\end{itemize}
The outputs are under \path{~/.cache/collisionAvoidance/ablation-20261005}.
\end{document}
